% =========================================================
% frontmatter/portfolio_context.tex
% =========================================================

\chapter*{Portfolio Context and Evidence}
\thispagestyle{fancy}
\addcontentsline{toc}{chapter}{Portfolio Context and Evidence}

\noindent
\textbf{Project Context.}
This work originated from a graduate-level antenna engineering simulation
completed during Summer 2026. The present document is an independently
prepared public portfolio edition intended to demonstrate guided-wave
analysis, CST full-wave simulation, multimode interpretation, numerical
verification, and technical documentation.

\vspace{0.5em}

\noindent
\textbf{Public Portfolio Edition.}
The original academic report has been reorganized and edited for
professional and educational use. Course-specific instructions, grading
material, student-identification information, institutional branding,
instructor information, and administrative content have been removed.
This document does not represent an official publication or endorsement
by any university, instructor, software vendor, or other institution.

\vspace{0.5em}

\noindent
\textbf{Evidence Classification.}
The results presented in this report are based on closed-form
rectangular-waveguide calculations, an executable Python verification,
and CST Studio Suite simulations. No waveguide hardware was fabricated,
and no VNA, field-probe, or antenna-range measurement campaign was
conducted. The reported cutoff frequencies, impedances, scattering
parameters, fields, and surface-current patterns are therefore classified
as analytical or simulated evidence, not measured performance.

\vspace{0.5em}

\noindent
\textbf{Simulation Scope.}
The model is an ideal vacuum-filled, PEC-walled, uniform square waveguide.
The public package includes selected CST figures but not the native binary
model or raw exported field and S-parameter datasets. Formal mesh,
port-basis, and calibration-plane convergence studies were not completed.
Finite conductivity, manufacturing tolerances, transitions, and a horn
flare remain outside the present scope and are discussed explicitly as
limitations.

\vfill

\begin{center}
    \textit{Prepared as part of the Hossein Molhem Engineering Portfolio.}
\end{center}
